\chapter{Building, linking and testing}
\label{ch:building}

\section{Requirements}

\begin{itemize}
  \item a C99 compiler (GCC or Clang are tested);
  \item CMake 3.13 or newer;
  \item an MPI implementation with its C bindings (the tests use MPICH's \texttt{mpirun});
  \item optionally LAPACK and BLAS, for the least-squares paraboloid fit of the remesher. Without
    LAPACK the fit falls back to a barycentric (butterfly) fit; the library is complete either
    way.
  \item optionally a \LaTeX{} installation (\texttt{pdflatex}, TikZ, \texttt{listings}) to build
    this manual. The library build never needs \LaTeX.
\end{itemize}

\section{Configuring and building}

\mftl{} is a CMake project that builds one static library, \texttt{libMFTL2.a}, from all
sources under \texttt{src/} (collected with \texttt{GLOB\_RECURSE}, so adding a source file
needs only a re-configure). A typical release build:
\begin{lstlisting}[language=bash,numbers=none]
cmake -S . -B build-release -DCMAKE_BUILD_TYPE=Release -DMFTL_USE_LAPACK=ON
cmake --build build-release -j
\end{lstlisting}
The options are listed in Table~\ref{tab:cmake}. When \mftl{} is configured as the top-level
project, the tests, the examples and (in Debug builds) the sanitizers are enabled by default;
when it is consumed through \texttt{add\_subdirectory} they are off.

\begin{table}[ht]
\centering\small
\begin{tabularx}{\textwidth}{@{}llX@{}}
\toprule
Option & Default & Meaning \\
\midrule
\texttt{CMAKE\_BUILD\_TYPE} & --- & \texttt{Release} (\texttt{-O3}, link-time optimisation) or \texttt{Debug} \\
\texttt{MFTL\_USE\_LAPACK} & OFF & link LAPACK/BLAS and define \texttt{MFTL\_HAVE\_LAPACK} (paraboloid fit with \texttt{dgelsy}) \\
\texttt{ENABLE\_SANITIZERS} & ON (top level) & in Debug builds with GCC/Clang: instrument library, tests and examples with ASan and UBSan; the archive then carries the sanitizer link flag to its consumers \\
\texttt{MFTL\_BUILD\_TESTS} & ON (top level) & build the test suite and register it with CTest \\
\texttt{BUILD\_EXAMPLES} & ON (top level) & build the programs in \texttt{examples/} \\
\texttt{MFTL\_ENABLE\_IPO} & ON & link-time optimisation of the Release library (GCC, fat LTO objects so non-LTO consumers still link) \\
\texttt{MFTL\_NATIVE\_ARCH} & OFF & \texttt{-march=native} for the Release library (not portable) \\
\texttt{MFTL\_BUILD\_MANUAL} & OFF & add the \texttt{manual} target that builds this manual \\
\bottomrule
\end{tabularx}
\caption{CMake options of \mftl.}
\label{tab:cmake}
\end{table}

The script \texttt{build\_all.py} wraps the common configurations: no argument builds a Release
tree in \texttt{build/}, \texttt{--debug} a Debug tree with sanitizers, \texttt{--valgrind} a
Debug tree without sanitizers, and \texttt{--lib-only} skips the examples.
\texttt{run\_local\_tests.py} builds and runs the test suite.

\section{Using MFTL2 from another project}

The library exports the CMake package \texttt{MFTL2} with the imported target
\texttt{MFTL2::MFTL2}. Its usage interface carries the public include directories (the headers
in \texttt{include/} and \texttt{khash.h}, which the public structures use), MPI, \texttt{libm}
and --- when the library was built with it --- LAPACK and the definition
\texttt{MFTL\_HAVE\_LAPACK}. A consumer names nothing but the target:
\begin{lstlisting}[language=bash,numbers=none]
# install MFTL2 (or point CMAKE_PREFIX_PATH at a build tree, which also exports the package)
cmake --install build-release --prefix $HOME/opt/mftl2
\end{lstlisting}
\begin{lstlisting}[numbers=none,language={}]
find_package(MFTL2 CONFIG REQUIRED)
target_link_libraries(my_solver PRIVATE MFTL2::MFTL2)
\end{lstlisting}
The directory \texttt{external\_link\_test/} is such a consumer; it is used to check an
installation. \multi{} consumes \mftl{} in exactly this way (its CMake variable
\texttt{MFTL2\_ROOT} is the search prefix).

A program includes the public headers it needs; Table~\ref{tab:headers} lists them.

\begin{table}[ht]
\centering\small
\begin{tabularx}{\textwidth}{@{}lX@{}}
\toprule
Header & Contents \\
\midrule
\texttt{mftl\_core.h} & data structures, object and ensemble lifecycle, per-triangle payload, sphere generator \\
\texttt{mftl\_constants.h} & return codes, tolerances, flags, enumerations, twin-state codes \\
\texttt{mftl\_topology.h} & half-edge navigation, twin-state predicates, one-ring traversal, builders, validation, ghost completion, boundary loops \\
\texttt{mftl\_communication.h} & ownership, redistribution, ghost refresh and ghost layers, MPI record types \\
\texttt{mftl\_geometry.h} & periodic primitives, triangle properties, integrals, local-frame kernels, transforms, generators, curvature, point queries \\
\texttt{mftl\_images.h} & image counters, canonicalisation, wrap flags, body reference point \\
\texttt{mftl\_phase\_set.h} & phase-boundary sets \\
\texttt{mftl\_io.h} & VTK/VTU output, GTS, native restart meshes \\
\texttt{mftl\_bvh.h} & LBVH, triangle intersection, spatial ghosts \\
\texttt{mftl\_overlap.h} & overlap volumes, 2-D convex hull and polygon helpers \\
\texttt{mftl\_remeshing.h} & remesh parameters, operators, smoothing, volume conservation, serial and parallel drivers, load balancing \\
\texttt{mftl\_memory.h}, \texttt{mftl\_vectoroperations.h}, \texttt{mftl\_arrayfunctions.h} & small helpers (growable arrays, vector algebra, integer sorting) \\
\bottomrule
\end{tabularx}
\caption{Public headers.}
\label{tab:headers}
\end{table}

\section{Source layout}
\label{sec:layout}

The sources are grouped by subsystem (Stage~17). Consumers never include anything from
\texttt{src/}; the public headers are in \texttt{include/} (\texttt{mftl\_boundary\_internal.h} is
an internal header that stays there because the test suite uses it).

\begin{center}\small
\begin{tabularx}{\textwidth}{@{}lX@{}}
\toprule
Directory & Contents \\
\midrule
\texttt{src/core/} & object lifecycle (\texttt{object.c}), half-edge topology and builders (\texttt{topology.c}, \texttt{soup.c}), twin states, validation, discovery and ghost completion (\texttt{twin\_states.c}), radix sort \\
\texttt{src/parallel/} & ensembles and decompositions (\texttt{ensemble.c}); MPI record types (\texttt{comm\_types.c}); redistribution (\texttt{redistribute.c}); ghost refresh (\texttt{ghost\_refresh.c}); ghost-layer extension, spatial ghosts and distance-based ghosting (\texttt{ghost\_layers.c}); private declarations in \texttt{comm\_internal.h} \\
\texttt{src/periodic/} & image counters, canonicalisation, wrap flags (\texttt{images.c}); phase-boundary sets (\texttt{phase\_set.c}) \\
\texttt{src/geometry/} & integrals (\texttt{geometry\_object.c}), local-frame kernels, curvature, transforms, point queries; private declarations in \texttt{geometry\_internal.h} \\
\texttt{src/intersection/} & LBVH and triangle intersection (\texttt{bvh.c}), overlap volumes, 2-D convex hull \\
\texttt{src/generators/} & sphere and the other generators \\
\texttt{src/io/} & VTK/VTU, GTS and native mesh I/O \\
\texttt{src/remeshing/} & operators and gates (\texttt{remeshing\_core.c}), smoothing and quality (\texttt{remesh\_smoothing.c}), the serial driver (\texttt{remesh\_driver.c}), vertex fitting, volume-conserving smoothing, the generalized boundary rules, \mftlone{} criteria, the parallel driver and load balancing; private declarations in \texttt{remeshing\_internal.h} \\
\texttt{src/internal/} & declarations shared by several source directories (twin-state helpers); a private include directory, never installed \\
\texttt{tests/}, \texttt{examples/} & the test suite and the example programs (including the programs of this manual) \\
\texttt{tools/} & \texttt{check\_twin\_gid\_readers.sh} (audit of raw twin-state reads) and the executable phase-set model \\
\bottomrule
\end{tabularx}
\end{center}


\section{Running the tests}

The test suite (129 CTest entries) covers every subsystem on 1, 2, 4 and 8 ranks, including
the Stage-16 acceptance tests (\texttt{Stage16A}--\texttt{Stage16F}) and the manual's example
programs (\texttt{Manual-<program>-np<N>}).
\begin{lstlisting}[language=bash,numbers=none]
cd build-release && ctest -j6 --output-on-failure
\end{lstlisting}
Notes:
\begin{itemize}
  \item The tests start \texttt{mpirun -np N} themselves. With MPICH, running more ranks than
    cores needs no flag (MPICH does not accept Open~MPI's \texttt{--oversubscribe}).
  \item In sanitizer builds, LeakSanitizer reports small allocations inside MPICH; the tests run
    with the suppression file \texttt{tests/lsan.supp}.
  \item Setting the environment variable \texttt{MFTL\_VALIDATE} to a non-zero value (\eg
    \texttt{MFTL\_VALIDATE=1}) makes Release builds run the
    full collective twin-state validation at every quiescent point (Debug builds always do,
    \S\ref{sec:validation}).
\end{itemize}

\section{Building this manual}

The manual lives in \texttt{doc/}: \texttt{MFTL2.tex}, the chapters in \texttt{doc/chapters/}
and the TikZ figures inline. Build it with
\begin{lstlisting}[language=bash,numbers=none]
make -C doc          # doc/MFTL2.pdf (pdflatex, three passes)
\end{lstlisting}
or configure with \texttt{-DMFTL\_BUILD\_MANUAL=ON} and build the target \texttt{manual}. The
listings are read from \texttt{examples/}, so the manual always shows the programs that the
tests compile and run.
